\chapter*{Abstract}

Large language models can learn a task from a handful of labelled examples placed in their prompt, a behaviour called in-context learning.
Which examples are shown matters, and a growing literature selects them by similarity, diversity or learned scores.
Little is known about whether selection helps when the test data shift away from the examples, or when what the examples teach contradicts what the model already believes.
Smaller models are known to follow such contradicting evidence only partially.

This thesis asks whether choosing the demonstrations can help a 7B model follow them over its own priors under controlled distribution shift.
We build a synthetic tabular testbed in which covariate shift, spurious-feature reversal and concept shift (realised by giving features names whose everyday meaning agrees or disagrees with the rule) can each be switched on alone, with known ground truth.
On Qwen2.5-7B-Instruct with eight demonstrations, we compare seven hand-designed selection strategies against a zero-shot baseline.
We then train a set-transformer selector (\sata), first on structural targets and then on the model's own measured priors (\satami), and look inside the model with the logit lens, linear probes and activation steering.
Every confirmatory test was written down before its data were seen.

\textbf{Feature names, not demonstrations, decide most of what the model predicts.}
Flipping the names costs 0.20 \AUROC (95\% CI [0.13, 0.28], 22 of 24 tasks), more than any shift of the data.
No hand-designed strategy overcomes it.
\textbf{A trained selector beats every strategy on a stand-in learner, but not on the instruction-tuned model}: none of the six pre-registered \sata contrasts, and none of the four \satami contrasts, favours the selector after Holm correction.
In an exploratory comparison, \sata does help the base model, whose priors pull less.
\textbf{Inside the model, the answer forms only in its last few layers, and that is where the names take over}: its zero-shot guess is decodable from layer 2, the demonstrated rule hardly at all, and a steering direction for ``follow the demonstrations'' moves its outputs without making them more correct.

Within the scale we could study (7B models, eight numeric demonstrations, synthetic tasks, a small compute budget), selection alone cannot route around a strong prior.
The testbed, the contrast families and the diagnosis give future work at larger scale a controlled place to test what can.
